% ============================================================
% CHARACTER REFERENCE C — Manager-Level Colleague (Uni Care, Inc.)
% Form LIC 301E
% Perspective: Mid-level manager / close day-to-day coworker —
% warm, specific, ground-level observations of conduct and growth.
% ============================================================
\newdocument{Character Reference --- Manager-Level Colleague}
\resetlicquestions

\licformheader

\licwriterinfo
  {\RefCName}
  {\RefCAddress}
  {\RefCCityStateZip}
  {\RefCPhone}

% ----------------------------------------------------------
\licquestion{How long have you known the person you are writing this reference for?}

I have known \ApplicantName\ for \RefCKnownDuration, since \RefCKnownSince, as colleagues at \EmployerName. Our working relationship has been close and consistent, involving frequent day-to-day interaction, and it has given me a clear and practical view of who he is as a professional and as a person.

% ----------------------------------------------------------
\licquestion{How do you know this person?}

\ApplicantName\ and I work together at \EmployerName, the parent organization of \EmployerDivisions. I serve as \RefCTitle, and \ApplicantName\ serves as the \EmployerRole. In the ordinary course of our work, our paths cross daily: he supports the systems and equipment my team depends on, responds to issues that surface across our field operations, and partners with me on the practical problems that arise when a home health, hospice, or home care organization serves elderly patients across a wide service area. Beyond our formal responsibilities, we have worked alongside each other long enough and closely enough that I have had the opportunity to observe his character not only in planned meetings, but in the small, unscripted moments where character is most visible---how he treats staff under stress, how he handles mistakes, how he speaks about patients when no one is watching.

% ----------------------------------------------------------
\licquestion{Please give your opinion of this person's character.}

\ApplicantName\ is one of the most dependable, humble, and genuinely kind people I have worked with. He is the colleague who picks up the phone when something breaks on a Saturday night, who follows up to make sure a small problem did not grow into a larger one, and who treats the front-desk staff with the same respect he gives to the executive team. He does his work quietly and thoroughly, and he takes responsibility for outcomes rather than deflecting to process or to other people. In a busy healthcare operation, that kind of consistency is rare and valuable.

I am aware that \ApplicantName\ has two misdemeanor DUI convictions, one from \IncidentOneDate\ and one from \IncidentTwoDate. He told me about them himself. He did not have to, and he did not soften the facts. He spoke plainly about what happened, took full responsibility, and described what he was doing about it. I remember that conversation clearly---not because it changed how I saw him, but because it confirmed something I had already come to believe about his character: that he would rather be known honestly than thought of well.

In the time since, I have watched \ApplicantName\ do the difficult, unglamorous work of personal change. He enrolled voluntarily at \RehabProgram\ for \RehabDuration---something no court required of him, and something he did at his own expense and on his own time. He kept showing up at work. He continued to meet his responsibilities without excuse. He made real, visible changes to his lifestyle: regular fitness, healthier habits, a more disciplined daily rhythm. He talks differently about himself now---more measured, more reflective, more aware of the impact his actions have on the people around him. The \ApplicantName\ I work with today is a stronger and steadier version of the person I first met, and the change did not come from a program certificate---it came from sustained, honest work on himself.

I trust him, and I would trust him with the care of a family member of my own.

% ----------------------------------------------------------
\licquestion{Please describe any interaction you have observed between this person and the client group he/she is requesting to work with. For example: Clients may be children, developmentally disabled children or adults, mentally impaired adults, or elderly.}

Because my role places me close to the day-to-day operations that serve our elderly patients, I have had many opportunities to observe \ApplicantName's interactions with this population and with the environments in which they are cared for:

\begin{itemize}[leftmargin=2em, itemsep=0.2em]
  \item I have accompanied \ApplicantName\ on on-site technology support calls at the homes of elderly home health and hospice patients. In every case, I have watched him enter the home with a level of care that I would want from anyone visiting an elderly loved one of mine: he removes his shoes when appropriate, speaks softly, greets the patient by name, explains what he is doing and why, and works efficiently so that the household returns to quiet as soon as possible. He does not treat the patient as an obstacle to his work; he treats the work as something that exists to serve the patient.
  \item When elderly patients or their family members have called our office for help with caregiver portals, remote monitoring devices, or scheduling systems, \ApplicantName\ has taken those calls personally on many occasions. I have overheard him walk an 80-year-old patient through a password reset, step by patient step, without a trace of impatience in his voice. That patience is not performed---it is who he is.
  \item In office conversations about specific patients, \ApplicantName\ speaks about them with dignity. He remembers details---the name of a spouse, a recent hospitalization, a caregiver's concern---that someone outside the clinical staff would not normally retain. When a patient we have served has passed away, he takes a moment; he does not simply move on.
  \item I have observed him step forward, without being asked, to help a family member or patient who seemed disoriented or confused in our office. Those are the moments that reveal a person's instincts, and his instincts are good.
\end{itemize}

The desire he has expressed to serve the elderly in a residential care setting is, in my observation, an accurate extension of who he already is.

% ----------------------------------------------------------
\licquestion{Please add any comments you feel are relevant about this person and his/her desire to work in a community care facility.}

I want to add the following observations for the Department's consideration:

\textbf{Honest Disclosure:} \ApplicantName\ chose to tell me about his record at a time when he did not have to. In my experience, that kind of unforced honesty is one of the most reliable windows into a person's character. He is someone who would rather face a difficult conversation than live with a hidden one, and that is the kind of person I believe the Department wants serving in a community care facility.

\textbf{Reliability Through Personal Difficulty:} I watched \ApplicantName\ go through the most difficult period of his life without ever letting it compromise the people who depended on him at work. He kept his commitments, met his deadlines, and supported his colleagues. That tells me everything I need to know about how he would show up for residents, families, and staff in a community care setting.

\textbf{Preparation for the Role:} \ApplicantName\ has not approached \FacilityTypeShort\ administration casually. He has completed the \RCFETraining. He has spent over five years inside a California-licensed healthcare organization where he has observed, up close, what excellent care looks like and what its operational requirements are. He holds \EducationDegrees, which reflect his capacity for rigorous, complex work. He is prepared.

\textbf{A Personal Recommendation:} I understand the weight of the recommendation I am giving. I am not vouching for \ApplicantName\ because of anything I have been asked to say; I am vouching for him because I know him, I have watched him, and I genuinely believe California's elderly residents will be safer and better cared for with him in this role than without him.

I wholeheartedly recommend \ApplicantName\ for the criminal record exemption and for \FacilityTypeShort\ Administrator Certification.

\licsignatureblock
